\anexoItem{Acta de Constitución del Proyecto}{
    \begingroup
    \renewcommand{\baselinestretch}{1}
    \footnotesize
    \setlength{\tabcolsep}{4pt}
    \renewcommand{\arraystretch}{1.2}

    % Primera tabla: Información General
    \begin{longtable}{@{}p{3.5cm}p{11.5cm}@{}}
        \caption{Acta de Constitución del Proyecto - Información General}
        \label{tab:anexo_acta_constitucion_general} \\
        \toprule
        \textbf{Componente} & \textbf{Detalle del Proyecto} \\ \midrule
        \endfirsthead
        
        \toprule
        \textbf{Componente} & \textbf{Detalle del Proyecto} \\ \midrule
        \endhead
        
        \midrule
        \multicolumn{2}{r}{\textit{(Continúa en la siguiente página)}} \\
        \endfoot
        
        \bottomrule
        \endlastfoot

        \textbf{Nombre del Proyecto} & 
        Desarrollo e implementación de un sistema web con flujos automatizados para la gestión de nóminas y procesos de recursos humanos orientado a clientes empresariales de J\&P Periféricos S.A.C. \\ \addlinespace

        \textbf{Breve descripción} & 
        Solución web modular construida con TypeScript, NestJS, Fastify, React 19 y PostgreSQL, integrada con el orquestador n8n, centraliza el legajo del personal, valida la asistencia biométrica, automatiza el cálculo masivo de planillas (horas extras, retenciones de ONP/AFP, quinta categoría, aportes de EsSalud) y genera de forma asíncrona archivos regulatorios (SUNAT PLAME, AFPnet, Telecrédito BCP) con boletas de pago electrónicas en PDF cifradas con el DNI del trabajador. \\ \addlinespace

        \textbf{Propósito del Proyecto} & 
        Erradicar la dependencia de hojas de cálculo descentralizadas y tareas manuales propensas a errores en las empresas clientes atendidas por J\&P Periféricos S.A.C., mitigando el riesgo de multas y contingencias ante SUNAT y SUNAFIL. Asimismo, dotar a la empresa patrocinadora de un activo tecnológico escalable para su comercialización bajo el modelo de suscripción en la nube (Software as a Service - SaaS). \\ \addlinespace

        \textbf{Objetivos del Proyecto} & 
        \begin{minipage}[t]{11.2cm}
        \begin{itemize}[nosep, leftmargin=*]
            \item Desarrollar una plataforma web centralizada para la administración de legajos, jornadas laborales y solicitudes del personal.
            \item Automatizar el cálculo masivo y periódico de nóminas aplicando las fórmulas laborales peruanas con precisión decimal exacta.
            \item Implementar flujos asíncronos con n8n para la generación de estructuras oficiales (PLAME, AFPnet, bancos) y distribución de boletas por correo.
            \item Garantizar la seguridad, trazabilidad e inmutabilidad de la información salarial mediante control de acceso por roles (RBAC) y bitácora de auditoría (audit log).
        \end{itemize}
        \end{minipage} \\ \addlinespace

        \textbf{Alcance del Proyecto} & 
        \begin{minipage}[t]{11.2cm}
        \begin{itemize}[nosep, leftmargin=*]
            \item Análisis de procesos y levantamiento formal de requerimientos (Scrum).
            \item Diseño de arquitectura de software (monolito modular) y modelo de datos relacional.
            \item Desarrollo del backend (NestJS/Fastify) y frontend interactivo (React 19 / Tailwind CSS / Shadcn).
            \item Integración de colas en memoria (BullMQ / Redis) y flujos M2M con n8n (firma HMAC-SHA256).
            \item Pruebas funcionales, unitarias (Jest), de integración y validación con usuarios clave.
            \item Despliegue en servidor físico local con Docker Compose y capacitación operativa.
        \end{itemize}
        \end{minipage} \\ \addlinespace

        \textbf{Justificación del Proyecto} & 
        J\&P Periféricos S.A.C. requiere modernizar y unificar la gestión laboral que ofrece a sus clientes empresariales, ventilando hojas de cálculo dispersas que generan demoras, reprocesos y riesgo de sanciones inspectivas. La implementación del sistema automatizado minimiza el tiempo de liquidación manual, garantiza la exactitud en retenciones de ley y amplía el portafolio comercial de la empresa mediante un servicio SaaS recurrente. \\ \addlinespace

        \textbf{Fecha de Emisión} & Martes, 12 de mayo de 2026 \\ \addlinespace
        \textbf{Patrocinador (Sponsor)} & Julio Florez Vicuña -- Gerente General de J\&P Periféricos S.A.C. \\ \addlinespace
        \textbf{Gerente del Proyecto} & Mg. José Armando Tisnado Utallán \\
    \end{longtable}

    \vspace{1em}

    % Segunda tabla: Equipo responsable
    \begin{longtable}{@{}p{1.2cm}p{5.0cm}p{6.8cm}c@{}}
        \caption{Equipo responsable del proyecto}
        \label{tab:anexo_acta_equipo} \\
        \toprule
        \textbf{Código} & \textbf{Nombre} & \textbf{Cargo} & \textbf{Horas} \\ \midrule
        \endfirsthead
        
        \toprule
        \textbf{Código} & \textbf{Nombre} & \textbf{Cargo} & \textbf{Horas} \\ \midrule
        \endhead
        
        \midrule
        \multicolumn{4}{r}{\textit{(Continúa en la siguiente página)}} \\
        \endfoot
        
        \bottomrule
        \endlastfoot

        PO1 & Julio Florez Vicuña & PRODUCT OWNER & 4 \\
        S1 & Juan Pablo Alcala & STAKEHOLDER & 2 \\
        S2 & Daniel Roman & STAKEHOLDER & 4 \\
        SC1 & Daniel Singer & SCRUM MASTER \& BACKEND LEAD & 8 \\
        T1 & Rodrigo Castillo & FRONTEND LEAD \& UI ENGINEER & 8 \\
        T2 & Dylan Florez & DEVOPS \& QA ENGINEER & 8 \\
        T3 & Adrian Dueñas & FULL STACK DEVELOPER & 8 \\
        T4 & Julio Vargas & JUNIOR DEVELOPER & 6 \\
        T5 & Luis Olazo Sisniegas & JUNIOR DEVELOPER & 6 \\
    \end{longtable}

    \vspace{1em}

    % Tercera tabla: Requerimientos a alto nivel (Márgenes corregidos)
    \begin{longtable}{@{}p{1.1cm}p{3.2cm}p{6.2cm}p{2.8cm}@{}}
        \caption{Requisitos del Proyecto a Alto Nivel}
        \label{tab:anexo_acta_requisitos} \\
        \toprule
        \textbf{Código} & \textbf{Nombre} & \textbf{Descripción} & \textbf{Rol} \\ \midrule
        \endfirsthead
        
        \toprule
        \textbf{Código} & \textbf{Nombre} & \textbf{Descripción} & \textbf{Rol} \\ \midrule
        \endhead
        
        \midrule
        \multicolumn{4}{r}{\textit{(Continúa en la siguiente página)}} \\
        \endfoot
        
        \bottomrule
        \endlastfoot

        RF-01 & Control accesos por roles (RBAC) & Autenticar usuarios con credenciales y emitir tokens JWT que restrinjan vistas y operaciones según el rol del usuario. & Personal de RRHH, Contador, Admin \\ \addlinespace
        RF-02 & Registrar eventos en bitácora de auditoría & Registrar automáticamente cada una de las críticas en base de datos capturando usuario, IP, tabla y diferencial (SOV). & Personal de RRHH, Contador, Admin \\ \addlinespace
        RF-03 & Administrar legajos digitales de empleados & Gestionar la información maestra del personal, régimen previsional, entidad bancaria, CCI y asignación familiar. & Personal de RRHH, Contador \\ \addlinespace
        RF-04 & Importar información laboral de forma masiva & Cargar lotes de colaboradores en dos pasos: prevalidación en memoria con reporte de errores y confirmación asíncrona. & Personal de RRHH, Asistente, Contador \\ \addlinespace
        RF-05 & Gestionar contratos y alertas de vencimiento & Administrar modalidades contractuales y alertas de forma visual de los vencimientos en rangos de 15, 30 y 60 días. & Personal de RRHH \\ \addlinespace
        RF-06 & Gestionar solicitudes y permisos del personal & Tramitar solicitudes con sustento digital (vacaciones, licencias) y determinar su efecto remunerativo en la asistencia. & Personal de RRHH, Empleado \\ \addlinespace
        RF-07 & Asignar y monitorear tareas operativas & Permitir al contador delegar tareas operativas al Asistente y auxiliar sus estados de revisión y aprobación. & Contador, Asistente Contable \\ \addlinespace
        RF-08 & Cargar archivos probatorios y notas en tareas & Permitir adjuntar sustentos digitales en PDF o imágenes y registrar anotaciones de seguimiento en las asignaciones. & Personal de RRHH, Asistente, Contador \\ \addlinespace
        RF-09 & Parametrizar variables normativas de ley & Actualizar valores vigentes de UIT, RMV, horas extras y comisiones de AFP preservando la vigencia histórica. & Contador, Asistente Contable \\ \addlinespace
        RF-10 & Distribuir masivamente boletas por correo & Integrar flujos en n8n para despachar vía SMTP los boletas de pago cifradas a las casillas de correo de cada trabajador. & Personal de RRHH, Asistente, Contador \\ \addlinespace
        RF-11 & Consolidar incidencias mensuales de asistencia & Importar marcaciones del reloj biométrico y cruzar faltas, tardanzas y sobretiempos con las solicitudes aprobadas. & Personal de RRHH, Asistente, Contador \\ \addlinespace
        RF-12 & Ejecutar cálculo masivo de nóminas & Procesar en hilos secundarios las remuneraciones brutas, retenciones (ONP/AFP, 5ta Cat.) y aportes de EsSalud con precisión decimal. & Contador, Asistente Contable \\ \addlinespace
        RF-13 & Ejecutar cierre inmutable y reapertura excepcional & Congelar la planilla aprobada mediante triggers SQL para impedir mutaciones, permitiendo reaperturas auditadas. & Contador \\ \addlinespace
        RF-14 & Generar boletas electrónicas en formato PDF & Renderizar automáticamente las boletas individuales en PDF a partir del snapshot inmutable de cálculo al cerrar el periodo. & Personal de RRHH, Asistente del Contador \\ \addlinespace
        RF-15 & Exportar archivos estructurados y soportes & Compilar estructuras oficiales para PDF PLAME (\texttt{.rem}, \texttt{.jor}), AFPnet, Telecrédito BCP y reportes de costos en CSV. & Contador, Asistente del Contador \\
    \end{longtable}

    \vspace{1em}

    % Cuarta tabla: Recursos Financieros y Presupuesto
    \begin{longtable}{@{}p{0.8cm}p{3.2cm}p{6.4cm}r@{}}
        \caption{Recursos Financieros y Presupuesto del Proyecto}
        \label{tab:anexo_acta_presupuesto} \\
        \toprule
        \textbf{N\textdegree} & \textbf{Descripción} & \textbf{Concepto} & \textbf{Monto} \\ \midrule
        \endfirsthead
        
        \toprule
        \textbf{N\textdegree} & \textbf{Descripción} & \textbf{Concepto} & \textbf{Monto} \\ \midrule
        \endhead
        
        \midrule
        \multicolumn{4}{r}{\textit{(Continúa en la siguiente página)}} \\
        \endfoot
        
        \bottomrule
        \endlastfoot

        1 & Personal & Equipo de Desarrollo Core (4 desarrolladores Mid-Level S/ 12,600 c/u) & S/ 50,400.00 \\
        2 & Personal & Equipo de Desarrollo Junior / Soporte (2 desarrolladores S/ 3,150 c/u) & S/ 6,300.00 \\
        3 & Material & Memoria RAM 32 GB DDR4 para Servidor Local & S/ 399.00 \\
        4 & Material & Unidad SSD Kingston A400 de 240 GB para Servidor Local & S/ 135.00 \\
        5 & Equipos / Software & Servicio de dominio y enrutamiento DynDNS Pro (1 año) & S/ 200.00 \\
        6 & Equipos / Software & Depreciación contable anual de 6 equipos de cómputo portátiles & S/ 2,253.66 \\
        7 & Otros Costos & Créditos de procesamiento API Inteligencia Artificial (Lora) & S/ 150.00 \\
        8 & Otros Costos & Servicios de internet y electricidad durante el desarrollo & S/ 750.00 \\
        9 & Otros Costos & Capacidad de interna en herramientas nuevas (si aplica) & S/ 625.00 \\ \midrule
        \multicolumn{3}{r}{\textbf{Total}} & \textbf{S/ 60,006.79} \\
    \end{longtable}

    \vspace{0.5cm}

    \begin{itemize}[leftmargin=*]
        \item \textbf{Cronograma Estimado del Proyecto:} Duración: 156 días laborables (12/05 al 29/12/2026).
        \item \textbf{Presupuesto Estimado del Proyecto:} S/ 60,006.79 (Aprox. \$17,705.45 USD).
    \end{itemize}

    \vspace{0.5cm}
    \textbf{Aprobación del Proyecto:} \\
    Se autoriza formalmente el inicio del proyecto de innovación y mejora. El equipo del proyecto cuenta con la autoridad técnica para aplicar los recursos asignados y coordinar con los representantes de las áreas de Recursos Humanos y Contabilidad de J\&P Periféricos S.A.C.
    \endgroup
}